% Ad Atticum 9.12 --- headnote
% ad-atticum-09-12, 20 March 49 BC, Formian villa

\textit{Cicero to Atticus, written from the Formian villa
on the thirteenth day before the Kalends of April 49 BC
(the manuscript dateline: \textit{Scr.\ in Formiano xiii
K.\ Apr.\ a.\ 705 (49)}), the same day as 9.11. A second
dispatch has overtaken the first: a letter from Lepta
reports that Pompey is shut in at Brundisium by a
siege-wall and that the harbour mouth is blocked by
rafts. The same news is independently confirmed by
Matius and Trebatius, whose couriers crossed Caesar's
at Minturnae. The report will turn out to be false ---
Pompey had already sailed --- but Cicero does not yet
know that, and the letter is written under the belief
that the man whose cause he is debating joining is on
the point of being taken. The tears are real, the
register breaks; ``so help me, for the tears, I cannot
think out the rest.''}

\textit{Section 2 turns aside for a flash of contempt
for Dionysius, the Greek freedman tutor of Cicero's
son and nephew, who has thrown over the household now
its fortune has fallen: not Panaetius to Cicero's
Scipio after all. Section 3 sets the unbearable
contrast: the army of the Roman people lays siege to
Pompey, and meanwhile in Rome the praetors give
judgement, the aediles get up the games, the honest
men enter up their interest payments --- ``and I myself
sit still.'' The letter closes on the wish for ``the
Mucian end'' (Q.\ Mucius Scaevola the pontifex, killed
in the Marian terror of 82), and the bleak last
sentence: nothing left to wish for, except to be set
free by some mercy of an enemy.}
